% =====================================================================
%  FICHE 8 — THÈME 3 — NIVEAU MOYEN — CORRIGÉ
% =====================================================================
\documentclass[11pt,a4paper]{article}
\usepackage[utf8]{inputenc}
\usepackage[T1]{fontenc}
\usepackage{lmodern}
\usepackage[french]{babel}
\usepackage{amsmath,amssymb,mathtools}
\usepackage{icomma}
\usepackage{siunitx}
\sisetup{output-decimal-marker={,},per-mode=power,group-separator={\,},group-minimum-digits=5}
\usepackage[version=4]{mhchem}
\usepackage{xcolor}
\usepackage{tikz}
\usepackage[most]{tcolorbox}
\usepackage{enumitem}
\usepackage{array}
\usepackage{fancyhdr}
\usepackage[hidelinks]{hyperref}
\usetikzlibrary{arrows.meta,positioning,calc}
\usepackage[a4paper,margin=2.1cm,headheight=26pt,headsep=10pt]{geometry}
\usepackage{titlesec}

\definecolor{primary}{RGB}{150,98,162}
\definecolor{primarydark}{RGB}{92,52,110}
\definecolor{excol}{RGB}{0,135,90}

\tcbset{boxbase/.style={enhanced,breakable,boxrule=0.9pt,arc=2.5pt,left=3mm,right=3mm,top=2.3mm,bottom=2.3mm,fonttitle=\bfseries,coltitle=white,toptitle=0.6mm,bottomtitle=0.6mm}}
\newtcolorbox[auto counter]{correction}[1][]{boxbase,colback=excol!4,colframe=excol,title={\textsc{Corrigé}~\thetcbcounter\ \ifx\relax#1\relax\relax\else\textmd{--- \itshape #1}\fi}}

\newcommand{\rep}[1]{\textcolor{primarydark}{\textbf{#1}}}

\pagestyle{fancy}
\fancyhf{}
\fancyhead[L]{\small\textcolor{primarydark}{\textbf{Fiche 8} — Thème 3 : Calorimétrie}}
\fancyhead[R]{\small\textcolor{primarydark}{\textsc{Moyen — corrigé}}}
\fancyfoot[C]{\small\thepage}
\renewcommand{\headrulewidth}{0.8pt}
\renewcommand{\headrule}{\hbox to\headwidth{\color{primary}\leaders\hrule height \headrulewidth\hfill}}
\setlength{\parindent}{0pt}\setlength{\parskip}{3pt}

\begin{document}

\begin{center}
\begin{tikzpicture}
\node[rectangle,rounded corners=7pt,fill=excol,text=white,minimum width=16cm,
      minimum height=1.1cm,align=center,inner sep=6pt]
  {{\large\bfseries Thème 3 — Calorimétrie}\\[1pt]{\small Niveau \textsc{Moyen} — corrigé détaillé}};
\end{tikzpicture}
\end{center}

\begin{correction}[chauffage en plusieurs étapes]
\begin{enumerate}[label=\alph*),leftmargin=2em,topsep=2pt,itemsep=2pt]
  \item Étape 1 : échauffement de la glace ($-20\to0\,\si{\celsius}$), $Q_1=m\,c_{\text{glace}}\Delta T$.
        Étape 2 : fusion à \SI{0}{\celsius}, $Q_2=m\,L_{\text{fus}}$. Étape 3 : échauffement de l'eau
        liquide ($0\to50\,\si{\celsius}$), $Q_3=m\,c_{\text{eau}}\Delta T$.
  \item $Q_1 = 0{,}10\times2{,}060\times20 = \SI{4.12}{\kilo\joule}$ ;\quad
        $Q_2 = 0{,}10\times334 = \SI{33.4}{\kilo\joule}$ ;\quad
        $Q_3 = 0{,}10\times4{,}185\times50 = \SI{20.925}{\kilo\joule}$.
  \item $Q = Q_1+Q_2+Q_3 = 4{,}12+33{,}4+20{,}925 = \rep{\SI{58.445}{\kilo\joule}}$.
\end{enumerate}
\end{correction}

\begin{correction}[neutralisation dans un calorimètre]
\begin{enumerate}[label=\alph*),leftmargin=2em,topsep=2pt,itemsep=2pt]
  \item Volume total $50+50=\SI{100}{\milli\litre}$, soit $m=\SI{0.100}{\kilogram}$ ;
        $\Delta T = 30-20 = \rep{\SI{10}{\kelvin}}$.
  \item $Q = m\,c_{\text{eau}}\,\Delta T = 0{,}100\times4{,}185\times10 = \rep{\SI{4.185}{\kilo\joule}} = \rep{\SI{4185}{\joule}}$.
  \item La température \emph{monte} : la réaction dégage de la chaleur, elle est \rep{exothermique}.
\end{enumerate}
\end{correction}

\begin{correction}[température finale d'un mélange]
\begin{enumerate}[label=\alph*),leftmargin=2em,topsep=2pt,itemsep=2pt]
  \item $\sum Q_i = 0$ : $\;m_1 c\,(T_{\text{f}}-80) + m_2 c\,(T_{\text{f}}-20) = 0$. En divisant par $c$ :
        \[ 0{,}200\,(T_{\text{f}}-80) + 0{,}300\,(T_{\text{f}}-20) = 0. \]
  \item $0{,}5\,T_{\text{f}} = 0{,}200\times80 + 0{,}300\times20 = 16 + 6 = 22$, d'où
        $T_{\text{f}} = \dfrac{22}{0{,}5} = \rep{\SI{44}{\celsius}}$. (Valeur comprise entre \SI{20}{} et
        \SI{80}{\celsius}, plus proche de \SI{20}{} car la masse froide est plus grande : cohérent.)
\end{enumerate}
\end{correction}

\begin{correction}[enthalpie molaire de dissolution]
\begin{enumerate}[label=\alph*),leftmargin=2em,topsep=2pt,itemsep=2pt]
  \item $Q = m\,c_{\text{eau}}\,\Delta T = 0{,}200\times4{,}185\times(24-20) = 0{,}200\times4{,}185\times4 = \rep{\SI{3.348}{\kilo\joule}}$.
  \item La température \emph{augmente} : la dissolution \emph{dégage} de la chaleur, elle est \rep{exothermique} ($\Delta H<0$).
  \item Le calorimètre étant de capacité négligeable, toute la chaleur cédée par le sel est captée par
        l'eau : $|Q_{\text{cédée}}| = \SI{3.348}{\kilo\joule}$. Par mole :
        \[ \Delta H = \frac{-\,3{,}348}{0{,}50} = \rep{\SI{-6.696}{\kilo\joule\per\mol}} \approx \SI{-6.7}{\kilo\joule\per\mol}. \]
\end{enumerate}
\end{correction}

\end{document}
